\section{Outlook}\label{sec:outlook}

\paragraph{What has been done.} The dimension is $d=4$ throughout.
\begin{itemize}
\item \emph{Part~I} treats the trivial background $U=1$ on the torus, with blocks centred at the points of the unit lattice.
\item \emph{Part~II} extends the results to backgrounds close to $1$ in a global gauge, by perturbation.
\item \emph{Part~III} replaces the global hypothesis by local regularity: on every cube of size of order $M$ there is a gauge in which the background is close to $1$. It provides:
  \begin{itemize}
  \item Dirichlet operators on unions of blocks;
  \item a local proof of coercivity by IMS localization;
  \item locality of all operators in the domain and in the background;
  \item Ba\l aban's random-walk expansion for the scalar Green function; for the gauge-field Green function it is completed in factorized form in Part~IX (Section~\ref{sec:nested-walks}).
  \end{itemize}
\item \emph{Part~IV} constructs Ba\l aban's symmetric nonlinear averaging in four dimensions, with centred blocks and group averages over the coordinate-ordered contours, the averaged axial gauge and the gauge-adapted averaging, and proves the axial-to-Landau chart. The linearized averaging is exactly covariant ($E=0$), it is covariant under the translations and the hyperoctahedral group $B_4$, and Parts~II and~III apply to it: the real statements under properties (i)--(iv) of Remark~\ref{rem:general-averaging}, the analytic ones with property (v) in addition.
\item \emph{Part~V} solves the variational problem relative to a regular reference configuration. It gives existence, uniqueness and minimality of the background field, second-order and H\"older bounds, the Euler--Lagrange equation, analyticity in the coarse field, exponential decay of the derivative, and the lattice symmetry of the solution.
\item \emph{Part~VI} carries out the induction over the number of steps. This proves Ba\l aban's variational theorem on the torus with constants independent of $k$, including its lattice symmetry (Theorem~\ref{thm:balaban}). As a consequence, every real configuration with small curvature and current is regular (Corollary~\ref{cor:regularity}). So the hypothesis of Part~III is a consequence of gauge-invariant conditions.
\item \emph{Part~VII} writes one small-field renormalization group step, with Ba\l aban's exponential gauge fixing, exactly as a Gaussian integral around the background field (Lemma~\ref{lem:exp-gauge}, Theorem~\ref{thm:step}). It proves the linear input: coercivity, uniformly in $k$, and decay of the fluctuation covariance (Theorem~\ref{thm:step-coercive}, Corollary~\ref{cor:covariance}), analyticity and tree decay of the coarse effective action, and local determinant densities. The step commutes with the translations and with $B_4$ (Proposition~\ref{prop:symmetry}, Lemma~\ref{lem:gf-step-invariance}).
\item \emph{Part~VIII} localizes the dependence on the background, including large changes, by Dirichlet approximants and telescoping:
  \begin{itemize}
  \item for the fluctuation covariance and the coarse Hessian (Theorem~\ref{thm:local-covariance}, Corollary~\ref{cor:large-changes});
  \item for the determinant, which becomes a constant plus normalized one-loop activities of second order in the curvature (Corollary~\ref{cor:one-loop});
  \item for the interaction and the composed activities (Theorem~\ref{thm:interaction-local}, Corollary~\ref{cor:composed}). This uses interior pointwise bounds with derivatives for the Dirichlet operators (Theorem~\ref{thm:ID}). They are proved by absorption at all scales inside a single cube, for a local elliptic system in which two scalar potentials of the slice projection are additional unknowns (Section~\ref{sec:interior}).
  \end{itemize}
  This localization is of radius type. On the main line it is superseded by Parts~IX and~X. From Part~VIII the main line keeps the interior bounds, the Dirichlet Hessians and local fluctuation operators, the local determinant densities, the graded domains and the quadratic bound for normalized activities (Lemma~\ref{lem:quadratic-bound}).
\item \emph{Part~IX} gives the volume-type localization by random walks with weakening parameters: for the covariance, its square root, the coarse Hessian and the determinant in $L^2$ norms (Section~\ref{sec:nested-walks}); interior estimates with $L^p$ data (Section~\ref{sec:interior-lp}); pointwise $C^{1,\beta}$ bounds for the weakened propagators through tier norms (Section~\ref{sec:tiers}); and the weakened background map (Theorem~\ref{thm:weak-bg}).
\item \emph{Part~X} carries out the small-field step in Dimock's scheme, and iterates it over all scales in Ba\l aban's four-dimensional organization:
  \begin{itemize}
  \item complex references, including Ba\l aban's averaged backgrounds (Proposition~\ref{prop:AB}, Proposition~\ref{prop:complex-ext});
  \item the weakened step quantities, including the gauge-fixing pieces, and the localization lemma (Section~\ref{sec:step-quantities});
  \item the cutoffs on the fluctuation field and the white noise, and the Mayer and polymer expansion; by parity there is no term of order $g$ (Section~\ref{sec:cluster});
  \item the step in polymer form (Theorem~\ref{thm:sf-step});
  \item the parts of the flow that do not depend on its four-dimensional organization: the small-field parameters, the polymer grid and the reblocking, the conditions on the parameters (Remark~\ref{rem:parameters}), the scales of the torus, the first step and its formula (Proposition~\ref{prop:first-step}, Lemma~\ref{lem:first-step-formula}), the vacuum energy (Proposition~\ref{prop:vacuum}) and the bookkeeping of the modified model (Remark~\ref{rem:sf-reference});
  \item the generations of activities, their symmetrization over the symmetries of their grid, and their coefficient $\beta$, read off from a period density that has no topological and no anisotropic part (Section~\ref{sec:generations}, Proposition~\ref{prop:period-iso});
  \item the extraction of the marginal part of an old generation at every later scale (Section~\ref{sec:extraction}, Proposition~\ref{prop:extraction});
  \item the old activities, bounded uniformly in the scale, and one step in generation form, with $g_{k+1}^{-2}=g_k^{-2}-\beta_{k+1}$ (Proposition~\ref{prop:old-activities}, Theorem~\ref{thm:gen-step});
  \item the flow over all scales on a fixed torus, as long as the couplings stay in $(0,\gamma]$, with the vacuum energy and a single counterterm (Theorem~\ref{thm:sf-flow}); the continuity in the couplings (Proposition~\ref{prop:flow-continuity}, Corollary~\ref{cor:bare-continuity}); and the bounds per volume on the partition function of the small-field model (Corollary~\ref{cor:sf-stability}).
  \end{itemize}
  Theorems~\ref{thm:K} and~\ref{thm:L} summarize the step, and Theorem~\ref{thm:N} the flow.
\item \emph{Part~XI} computes the one-loop coefficient. It gives the relation between the one-loop terms of the composite maps of $k$ and $k-1$ steps (Lemma~\ref{lem:ol-one-step}) and, with the formula for the first step (Lemma~\ref{lem:first-step-formula}), the composition of the one-loop terms of the steps (Proposition~\ref{prop:ol-composition}); the formula of the composite one-loop term in the Landau chart (Theorem~\ref{thm:ol-chart}); its exact invariance under the translations and $B_4$, and that of the one-loop terms of the steps (Lemma~\ref{lem:ol-symmetry}); and from this invariance, with the infinite-volume limits of the second variations (Lemmas~\ref{lem:ol-torus-limit} and~\ref{lem:ol-limits}), a single coefficient $\beta$ and no term linear in $p$. The operators at $U=1$ are constructed in infinite volume by folding the torus operators (Section~\ref{subsec:ol-infinite-volume}), and the one-loop pieces of the polymers that do not wrap, and the Jacobian terms, are the same functions on all admissible tori (Proposition~\ref{prop:one-loop-local}). It also gives the logarithm on the fine lattice (Theorem~\ref{thm:ol-log}) and the cumulative coefficient $k\beta_\infty+O(1)$ (Theorem~\ref{thm:ol-cumulative}). With the flow of Part~X it gives the properties (F1)--(F3) of the couplings (Proposition~\ref{prop:flow-F}) and the ultraviolet limit, with the bounds on $g_k^{-2}-g_K^{-2}$ at every scale $k$ (Theorem~\ref{thm:ol-uv}). The comparison of the coefficients $\beta_j$, at fixed couplings and as the torus grows, with the coefficient that Ba\l aban's definition \cite[(1.20)--(1.22)]{BalYM4I} assigns to the generations (Proposition~\ref{prop:beta-identification}) rests on the generation locality of Section~\ref{sec:gen-locality}: the raw generations built with fixed couplings converge at the flat configurations as the torus grows, and so do their coefficients. Theorem~\ref{thm:M} summarizes Part~XI.
\end{itemize}

\paragraph{How the induction closes.} The reference at step $k$ is the solution of the $(k-1)$-step problem for the trivial extension of $V$, and it has exactly the prescribed averages. Theorem~\ref{thm:variational} then solves the problem at step $k$. Its regularity bound is not good enough to continue, because the change of units from $\delta_{k-1}$ to $\delta_k$ costs a power of $L$. The local improvement (Lemma~\ref{lem:improvement}) uses only Theorem~\ref{thm:variational} itself, applied with a different reference. This reference is spliced to the trivial configuration near a given point and coincides with the solution far away. Near the point, the solution is then $\bH$ applied to the coarse field in the gauge of the coordinate-ordered contours from the point, up to an error that is exponentially small in the splicing radius times the a priori bound, plus a quadratic error. The coarse field in this gauge is controlled by the curvature of $V$. Two features make this work on the torus:
\begin{itemize}
\item the uniformity of Theorem~\ref{thm:variational} in the reference;
\item the polynomial loss in the patching of local gauges (Lemma~\ref{lem:large-cubes}), which is beaten by the exponential decay of $\bH$ and $\bC$.
\end{itemize}

\paragraph{Why the global hypothesis was not enough.} Applied up to a real gauge transformation, hypothesis (B) is a condition on the real gauge orbit, so it does not break gauge invariance. But it is not local, and it is restrictive. By the isoperimetric argument of Section~\ref{sec:background}, every Wilson loop obeys a perimeter bound. Field strengths of size $\eps$ that are coherent over many blocks, and backgrounds with non-zero flux through a $2$-cycle, are excluded. Moreover the renormalization group needs small-field conditions that factorize over local regions, and operators whose dependence on distant changes of the background is weak. Part~III provides both: regularity is a local condition, and Theorem~\ref{thm:reg-locality} gives the locality. Part~VI shows that the local condition is implied by Ba\l aban's gauge-invariant small-field conditions.

A naive comparison of $G(U)$ and $G(U')$ through the resolvent identity would produce the difference $D_U-D_{U'}$, whose norm can be of order $\eta^{-1}$ where the backgrounds differ. The proof of Theorem~\ref{thm:reg-locality} avoids this. It compares each background with a Dirichlet operator on the region where they agree. Changing the domain at a fixed background produces cut couplings and differences of domain-dependent ingredients, including mass profiles and non-local terms, which are supported near the cut or exponentially small away from it. Derivative factors may occur at one or both ends of these terms; the estimates use the corresponding bounded undifferentiated or differentiated propagators.

\paragraph{Relation to Ba\l aban's construction.} Ba\l aban works with Dirichlet operators on domains \cite{BalBackground} and with local regularity on cubes \cite[(3.35)--(3.36)]{BalBackground}, and builds the global operators by random walks. His positivity comes from \cite[Theorem~3.11]{BalBackground} and local axial-gauge estimates \cite[Lemma~2.4]{BalPropII}. Our route differs in several technical points:
\begin{itemize}
\item coercivity comes from IMS localization together with the coercivity of Part~I;
\item the local gauge-field propagators carry a mass term near the boundary of their domain;
\item the normalization of the gauge transformation from the averaged axial gauge to the Landau gauge is built into a recursive correspondence map (Theorem~\ref{thm:correspondence}). Ba\l aban's nonlinear averages of gauge transformations then appear only through the identity for pure gauges;
\item the second-order regularity of the minimizer is derived from the Euler--Lagrange equation, which follows from gauge invariance and a transversality argument. We do not use the alternative equation \cite[(143)]{BalVar};
\item the local improvement of the regularity constants is obtained on the torus, by changing the reference configuration, instead of on nested domains (Lemma~\ref{lem:improvement}). The local regularity of general small-field configurations is derived from that of the background fields, through the chart (Corollary~\ref{cor:regularity}).
\end{itemize}
In Parts~I--VIII the random-walk expansion is used only for the explicit localized representation, not for the estimates. In Part~IX its weakened form (Section~\ref{sec:nested-walks}) also gives the bounds on the weakened operators, uniformly for $s\in\mathbb D_\kappa$, and their analyticity in the interior of $\mathbb D_\kappa$ (Theorem~\ref{thm:weak-inverse}, Corollary~\ref{cor:weak-covariance}, Lemma~\ref{lem:weak-basic}(a)). The polymer expansion of Part~X uses these bounds; the bounds on the derivatives in the weakening parameters come from Cauchy's estimate in the proof of Lemma~\ref{lem:localization}.

For Parts~VII--XI, Section~\ref{g:comparison} of the guide compares the note with Ba\l aban's and Dimock's papers item by item (Table~\ref{tab:G-comparison}). The main differences are these.
\begin{itemize}
\item The generations and the coefficients $\beta_j$ depend on the whole sequence of the earlier couplings, and only their continuity in the couplings is proved; Ba\l aban's are smooth functions of a single coupling (Remark~\ref{rem:sf-flow}(a), Proposition~\ref{prop:flow-continuity}).
\item The small-field threshold $h_k=p(g_k)g_k$ and the white-noise bound $p_0(g_k)$ depend on the coupling, as in \cite{DimockI}, so the small-field domain shrinks with the coupling (Remark~\ref{rem:sf-flow}(a)).
\item The marginal part of an old generation is extracted in a complex gauge chosen for each polymer, to second order, with a power gain for the polymers that are not expanded; Ba\l aban expands in a gauge common to a cube, to fifth order, with Ward identities (Section~\ref{sec:extraction}).
\item The later-scale bounds, which Ba\l aban's Theorem~1 asserts, are proved in the induction (Theorem~\ref{thm:sf-flow}(d), Remark~\ref{rem:sf-flow}(a)).
\item The ultraviolet behaviour of the couplings of the small-field flow is proved for every compact, connected and simple group, while Ba\l aban states it for $SU(2)$. For the last $m_1$ steps its lower half holds up to an additive constant (Theorem~\ref{thm:ol-uv} and the paragraph after it).
\end{itemize}

\paragraph{What is not proved.} The flow results of Parts~VII--X and Theorem~\ref{thm:M}(iv) concern the small-field model on a fixed torus (Remark~\ref{rem:sf-reference}(c)); Part~XI also uses infinite-volume limits of the one-loop terms (Theorem~\ref{thm:M}(ii),~(iii)). The large-field stage is not treated (Remark~\ref{rem:sf-reference}(d)), so nothing is proved about the partition function or the correlation functions of the lattice theory itself. No continuum limit is taken: Theorem~\ref{thm:M}(iv) concerns the couplings of the small-field flow, and no limit of the measure or of correlation functions as the lattice spacing tends to $0$ is considered. The other omissions are technical: a rate in Proposition~\ref{prop:beta-identification}, the exact lower half of the analogue of Ba\l aban's (0.31) for the last steps, linear bounds for the gauges on large cubes, gauge groups with abelian factors, and dimensions other than $4$ (Section~\ref{ss:1.5}).

\paragraph{Next steps.}
\begin{enumerate}[label=(\arabic*)]
\item \emph{A check by a human expert.} The note has not been checked by a human expert. The companion volume expands the proofs of Parts~VII--XI step by step, and its chapter on the parameter ledger records the constants and the order in which they are chosen. Section~\ref{g:hard} of the guide names the places to check first.
\item \emph{What remains of Part~XI} (Remark~\ref{rem:ol-remains}(c)). Generation locality is proved without a rate (Section~\ref{sec:gen-locality}). A rate, uniform in the couplings, is expected to follow from a Lipschitz bound of the step on the infinite lattice in the old activities, a bound on the errors of the old activities, and a condition on $\gamma$ (Remark~\ref{rem:gl-rate}(a)). For the exact lower half of the analogue of (0.31) at the last steps, a possible route is the convergence of the single one-loop coefficients, through a continuum limit at $U=1$ of the remainder of Proposition~\ref{prop:ol-reduction}, together with conditions on $n_K$ and $\gamma$. A smooth small-field cutoff, as in \cite[(2.9)]{BalYM4I}, should improve (F3) to $O(g^2)$ (Remark~\ref{rem:ol-remains}(a)).
\item \emph{Large fields and nested domains,} where Dirichlet operators on irregular domains are essential. This is the analogue of \cite{DimockII,DimockIII}. It must control what the small-field model discards at every step (Remark~\ref{rem:sf-reference}), and so remove the modification. Since the covariance $C_k$ depends on the background everywhere, the white-noise cutoff of Section~\ref{sec:cluster} will need a version that depends on regions. Part~III treats arbitrary unions of blocks, so this is the natural setting for them. The splicing of Part~VI is a torus device. On nested domains the reference no longer has exactly the prescribed averages, and the analysis on domains of Part~III must be extended to the nonlinear problem.
\item \emph{Gauges on cubes of size up to $c\,\eps_1^{-1}$} with bounds linear in the size (Remark~\ref{rem:large-cubes}).
\item \emph{Gauge groups with abelian factors.} The quadratic bound for normalized activities uses semisimplicity (Lemma~\ref{lem:quadratic-bound}). For abelian factors, gauge-invariant terms linear in the curvature can occur, and they must be treated separately.
\end{enumerate}
